\section*{Course outline}
\addcontentsline{toc}{section}{Course outline}

The course is planned along the following chapters. Chapters that are not yet covered by
these notes are listed for orientation only.

\begin{enumerate}[leftmargin=*, itemsep=0.6ex]
    \item \textbf{Special relativity (SR)}
          \begin{itemize}[nosep]
              \item Principles and spacetime diagrams
              \item Lorentz transformations
              \item Metric
              \item Geodesics
          \end{itemize}
    \item \textbf{SR in arbitrary coordinates} (``gravity'' from coordinates)
          \begin{itemize}[nosep]
              \item Equivalence principle
              \item $\Gamma$ from coordinate transformations
              \item Geodesic equation and metric
              \item Metric as gravitational potential: weak-field limit, gravitational redshift
          \end{itemize}
    \item \textbf{Curvature}
          \begin{itemize}[nosep]
              \item Equivalence principle $\Rightarrow$ tides are the true measure of gravity
              \item Tides $=$ curvature of spacetime
              \item Vectors and tensors: transformation properties
              \item Manifolds, differential structure
              \item Parallel transport, hence geodesics
              \item Geodesic deviation and curvature
              \item Connection to the weak-field limit, equivalence principle and gravitational redshift
          \end{itemize}
    \item \textbf{Vacuum solutions} (``pre'' Einstein equations)
          \begin{itemize}[nosep]
              \item Schwarzschild solution: geodesics, lensing, horizons
              \item Wave solutions: gravitational waves
          \end{itemize}
    \item \textbf{Matter}
          \begin{itemize}[nosep]
              \item Fluids in SR: $T_{\mu\nu}$
              \item Fluids in GR: $T_{\mu\nu}$
              \item Equation of motion: $\nabla_\mu T^{\mu\nu} = 0$
          \end{itemize}
    \item \textbf{Einstein equations}
          \begin{itemize}[nosep]
              \item $3+1$ split
              \item Application: FLRW (symmetry arguments)
          \end{itemize}
    \item \textbf{Cosmology}
\end{enumerate}

\newpage
